% !PRINTED 516
\weberpage{516}
%<<CONT
3. \gesperrt{Jedes der beiden Vorzeichen (11) kommt in gleich vielen Klassen der Diskriminante} $D=-4m'm''$ \gesperrt{vor und die Klassenfunktion ist durch Adjunktion von}
\[
\begin{array}{ll}
\sqrt{m''}, & \text{wenn } m''\equiv 1 \\
\sqrt{m''}, & ,,\quad m''\equiv 3
\end{array}
\ (\mathrm{mod}\ 4)
\]
\gesperrt{in zwei Faktoren vom Grade} $\tfrac12 h$ \gesperrt{zerlegbar, außer wenn} $m''\equiv 3$ \gesperrt{und} $m'$ \gesperrt{ein Quadrat ist.}
% SOURCE-ORIGINAL-SIC: The registered page repeats sqrt(m'') in both rows of the preceding case display; the m'' congruent to 3 case distinction is therefore mathematically suspect, but the printed reading is preserved.

In diesem Ausnahmefall ist, wenn wir $Q^2=4m'$ setzen, $-m''=\varDelta$ der Stamm von
\[
D = Q^2\varDelta,
\]
und es ist für jede durch eine Form der Diskriminante $D$ darstellbare und zu $D$ teilerfremde Zahl $A$:
\[
(D,A) = \left(\frac{A}{m''}\right) = +1
\]
und $H(u)$ ist mit $\Phi(M,u)$ identisch.

Zerlegt man $m$ in einer zweiten Art in zwei Faktoren
\weq{15}{m = m_1' m_1'',}
wo $m_1''$ denselben Bedingungen genügt wie $m''$, so erhält man in gleicher Weise eine Zerlegung
\weq{16}{H(u) = \Phi(M_1,u)\,\Phi(-M_1,u),}
und indem man den größten gemeinschaftlichen Teiler von $\Phi(M,u)$ und $\Phi(M_1,u)$ aufsucht, erhält man eine Zerlegung von $H(u)$ in vier Faktoren:
\[
H(u) = \Phi_1(u)\,\Phi_2(u)\,\Phi_3(u)\,\Phi_4(u),
\]
vorausgesetzt, daß in $\Phi(M,u)$ eine Quadratwurzel einer Primzahl vorkommt, die in $\Phi(M_1,u)$ nicht enthalten ist. So fährt man fort und zerlegt allmählich $H(u)$ in Faktoren, deren Anzahl eine Potenz von $2$ ist.

Wir wollen die Anwendung auf die einzelnen Fälle etwas genauer betrachten.

1) Ist $m\equiv 3 \ (\mathrm{mod}\ 4)$, so sind alle Charaktere in der Form
\[
\left(\frac{A}{m''}\right)
\]
enthalten, und die Formel (8) reicht hin, um \gesperrt{alle} Geschlechter voneinander zu trennen. Ist hier $m''\equiv 3$, so ist $m'\equiv 1 \ (\mathrm{mod}\ 4)$ und aus (12) folgt, daß in den Teilgleichungen nur die Quadratwurzeln aus solchen Zahlen vorkommen, die von der Form $4n+1$
%CONT>>
